\sectiontheme{rustaccent}
\section{Why Emi Exists}
\label{sec:overview}

Every team that ships an API ends up writing the same thing twice --- once
on the backend, where the route, the request body, and the response shape
live, and again on the frontend, hand-copied into a \texttt{fetch} call, a
model struct, a few \texttt{interface} declarations. Nobody decides to do
this; it just happens, and the contract between backend and frontend
isn't written down anywhere a compiler can see it. \textbf{Emi exists to
make that contract a real, compilable thing}, and then generate every side
of it, in idiomatic code, for every language a team ships.

\subsection{Where Emi sits relative to what you've tried}
\begin{itemize}[leftmargin=1.1em,itemsep=2pt,topsep=2pt]
  \item \textbf{vs.\ OpenAPI generators} --- an OpenAPI spec is downstream
        of the real API, a second artifact that has to be kept in sync by
        hand or another generator. Emi inverts this: the Emi definition
        \emph{is} the source, and your Go backend is generated \emph{from}
        it, so the spec and the server can never disagree. Emi can still
        emit a standard OpenAPI 3 document and a Postman collection on
        demand (\S\ref{sec:exports}) --- interoperability without making a
        hand-maintained spec the source of truth.
  \item \textbf{vs.\ gRPC/Protobuf} --- gRPC solves a real version of this
        problem but asks for a binary wire protocol browsers can't speak
        natively, a runtime dependency in every client, and an all-in
        ecosystem. Emi generates \textbf{plain HTTP and plain JSON}, with
        \textbf{no Emi runtime} in the output --- nothing to import, no
        library version to keep in lockstep. You can read every line Emi
        produces, hand-edit it, and eject at any time.
\end{itemize}

\begin{pullquote}
Emi's own one-sentence pitch: ``OpenAPI describes an API that already
exists. gRPC replaces your transport with its own. Emi generates the API
itself --- backend and every client --- as ordinary, runtime-free,
idiomatic code you fully own.''
\end{pullquote}

\subsection{How Emi generates code}
The whole system rests on one idea: \textbf{the definition is
target-agnostic}. An Emi definition (\texttt{.emi} YAML or JSON) describes
\emph{intent} --- actions, DTOs, enums, fields, nullability, collections
--- and says nothing about Go, TypeScript, or Swift. A series of
sub-compilers each turn that one definition into code for one target
(the full fan-out is diagrammed in Fig.~1, \S\ref{sec:go-entities}).
The CLI mirrors this directly:
\begin{lstlisting}
emi go      --path module.emi --output ./server
emi js      --path module.emi --output ./web/sdk \
  --tags react,typescript
emi swift   --path module.emi --output ./ios
emi kotlin  --path module.emi --output ./android
emi openapi --path module.emi
\end{lstlisting}
\texttt{--tags} let one compiler shape its output without touching the
definition: \texttt{--tags react} adds TanStack Query hooks plus
\texttt{useSse}/\texttt{useWebSocketX}; \texttt{--tags nestjs} emits
decorators that drop generated classes straight into Nest.js handlers;
\texttt{--tags typescript} switches the JS compiler from runtime JS to
\texttt{.ts}. None of them are mentioned in the definition --- they're
presentation choices made at compile time. Because the compiler is
written in Go and also compiled to WebAssembly, the \emph{exact same
compiler} runs in the browser --- that's what powers the
\href{https://torabian.github.io/emi/playground}{live playground}.

\subsection{Why classes, not interfaces}
Most codegen tools emit TypeScript \texttt{interface}s --- a compile-time
promise, erased at runtime, so a server that starts sending \texttt{null}
where it always sent a string throws in production while TypeScript
shrugged at compile time. Emi's JS compiler emits real \textbf{classes}
instead: private fields with typed getters/setters that coerce or reject
bad values, \textbf{auto-initialization} of every non-nullable field
(including nested objects, arrays and child DTOs, see
\S\ref{sec:js-compiler}), and real class instances handed back from the
network (built through a \texttt{creatorFn}), not a raw JSON blob cast to
a type. TypeScript checks types at compile time \emph{and} the class
enforces them at runtime --- the safety Go programmers take for granted,
in a language that historically had none of it.
